\chapter{Implementacja aplikacji}
\label{cha:implementacja}

\section{Najważniejsze moduły systemu}
\label{sec:moduly}

Logika aplikacji skupiona jest w warstwie usług. Kluczowe klasy to:

\begin{itemize}
  \item \texttt{DbRepository} - centralne repozytorium operacji na bazie danych; pobieranie
        i zapisywanie testów, grup, podejść i wyników oraz mapowanie encji na obiekty DTO,
  \item \texttt{GradingService} - logika oceniania (punkty za pojedyncze pytanie oraz przeliczanie
        całego wyniku podejścia),
  \item \texttt{CsvExporter} - eksport wyników testu do pliku CSV,
  \item \texttt{PasswordHash} - haszowanie haseł,
  \item \texttt{Safe} - wspólna obsługa wyjątków.
\end{itemize}

Wzajemne powiązania najważniejszych klas modelu przedstawia diagram UML na rysunku~\ref{rys:uml}.

\figscreen{uml.png}{0.95}{Diagram UML klas modelu aplikacji.}{rys:uml}

\section{Mechanizm logowania i autoryzacji}
\label{sec:logowanie}

Hasła nie są przechowywane jawnie - zapisywany jest ich skrót SHA-256 w postaci szesnastkowej.
Podczas logowania wyznaczany jest skrót wpisanego hasła i porównywany ze skrótem w bazie
(listing~\ref{lst:hash}). Autoryzacja oparta jest na roli użytkownika: po zalogowaniu obiekt
użytkownika trafia do okna głównego, które buduje inny zestaw modułów dla nauczyciela, a inny dla
ucznia.

\begin{lstlisting}[caption={Haszowanie hasła i logowanie.}, label=lst:hash]
public static string Hash(string password)
{
    byte[] bytes = SHA256.HashData(Encoding.UTF8.GetBytes(password));
    return Convert.ToHexString(bytes).ToLowerInvariant();
}

// logowanie - porownanie skrotu wpisanego hasla ze skrotem w bazie
var hash = PasswordHash.Hash(password);
var user = ctx.Users.FirstOrDefault(
    x => x.Username == username && x.Password == hash);
\end{lstlisting}

\section{Realizacja operacji CRUD}
\label{sec:crud}

Operacje tworzenia, odczytu, modyfikacji i usuwania (CRUD) realizowane są przez repozytorium
\texttt{DbRepository} z użyciem Entity Framework Core. Przykładem nietrywialnej operacji jest
usuwanie testu (listing~\ref{lst:delete}). Ponieważ odpowiedzi uczniów są powiązane z pytaniami
kluczem obcym w trybie ograniczającym (\texttt{Restrict}), bezpośrednie usunięcie testu mogłoby
naruszyć integralność. Dlatego najpierw usuwane są odpowiedzi i podejścia, a dopiero potem sam test -
pozostałe dane (pytania, opcje, przypisania grup) znikają kaskadowo.

\begin{lstlisting}[caption={Bezpieczne usuwanie testu (zachowanie integralności).}, label=lst:delete]
public static void DeleteTest(int id)
{
    using var ctx = Ctx();
    var quiz = ctx.Quizzes
        .Include(q => q.Attempts).ThenInclude(a => a.Answers)
        .FirstOrDefault(q => q.Id == id);
    if (quiz == null) return;

    foreach (var a in quiz.Attempts)
        ctx.StudentAnswers.RemoveRange(a.Answers);
    ctx.QuizAttempts.RemoveRange(quiz.Attempts);
    ctx.Quizzes.Remove(quiz);
    ctx.SaveChanges();
}
\end{lstlisting}

\section{Logika oceniania}
\label{sec:logikaOceniania}

Najważniejszym elementem logiki jest serwis oceniania. Metoda \texttt{Grade} wyznacza liczbę punktów za pojedyncze
pytanie (listing~\ref{lst:grade}). Pierwszeństwo ma ręcznie przyznana ocena - dotyczy to każdego
typu pytania, dzięki czemu nauczyciel może nadpisać również automatyczną ocenę pytania zamkniętego.
W przypadku braku ręcznej oceny pytania zamknięte są liczone automatycznie, przy czym w pytaniach
wielokrotnego wyboru przyznawane są punkty częściowe (proporcjonalnie do liczby trafień
pomniejszonych o błędne wskazania). Pytania otwarte bez ręcznej oceny oczekują na nauczyciela.
Cały wynik podejścia przeliczany jest następnie metodą \texttt{Recompute}, a przy wielu podejściach
do oceny brany jest najlepszy wynik ucznia.

\begin{lstlisting}[caption={Ocenianie pojedynczego pytania (GradingService).}, label=lst:grade]
public static (int pts, bool ok, bool partial, bool pending) Grade(
    Question q, object? ans, IReadOnlyDictionary<int, int>? manual = null)
{
    // reczne nadpisanie punktow ma pierwszenstwo (kazdy typ pytania)
    if (manual != null && manual.TryGetValue(q.Id, out var mo))
        return (mo, mo >= q.Points, mo > 0 && mo < q.Points, false);

    switch (q.Type)
    {
        case QuestionType.Multi:
            var chosen = AsIndices(ans);
            int good = chosen.Count(i => q.Correct.Contains(i));
            int bad  = chosen.Count(i => !q.Correct.Contains(i));
            double raw = q.Correct.Count == 0 ? 0
                         : Math.Max(0, good - bad) / (double)q.Correct.Count;
            int pts = (int)Math.Round(raw * q.Points);
            return (pts, pts == q.Points && q.Points > 0,
                    pts > 0 && pts < q.Points, false);
        // ... pytania jednokrotnego wyboru i otwarte
    }
    return (0, false, false, false);
}
\end{lstlisting}

\section{Organizacja nawigacji}
\label{sec:nawigacja}

Główne okno (\texttt{MainWindow}) pełni rolę powłoki: zawiera stały panel boczny z nawigacją oraz
obszar treści, w którym wyświetlane są kolejne ekrany modułów. Po kliknięciu pozycji menu tworzona
jest odpowiednia strona i wstawiana do obszaru treści. Ekran rozwiązywania testu oraz ekran wyniku
wyświetlane są w trybie pełnoekranowym (bez panelu bocznego), aby skupić uwagę ucznia na zadaniu.

\section{Interfejs użytkownika i stylizacja}
\label{sec:stylizacja}

Wygląd aplikacji zdefiniowano w słowniku zasobów \texttt{Theme.xaml}, obejmującym paletę kolorów oraz
style przycisków, pól i innych kontrolek. Powtarzalne elementy interfejsu (karty, znaczniki statusu,
wykresy kołowe i słupkowe, awatary) budowane są przez wspólne metody pomocnicze, dzięki czemu ekrany
wyglądają jednolicie. Okno główne otwiera się zmaksymalizowane, a układ oparty na siatkach dopasowuje
się do rozmiaru okna.
